\section{Redundancy Value}

Redundancy value is the annual freight cost reduction achievable by adding an adequate
alternate route to the network. It answers a different question than the closure cost
calculation: not ``how much does this corridor cost when it closes?'' but rather ``how
much of that cost could be eliminated by giving trucks a better alternate?''

\subsection{Donner Pass: Redundancy Value \$1.9 Billion per Year}

Current state: Donner Pass trucks reroute via I-40 (Flagstaff, Arizona) at a 550-mile
detour. This is the only interstate-quality alternate for a northern Sierra Nevada crossing.
The I-40 detour adds approximately 9.2 hours of driving time and crosses desert terrain with
limited truck services, generating the large B1 penalty multiplier (6.50) used in the
closure cost computation.

\textbf{Without any alternate}: if the I-40 alternate were not available (or were also
closed), trucks would face stranded freight conditions --- waiting indefinitely at the
closure point or diverting non-interstate routes (US-395 south, SR-89 over Monitor Pass)
that cannot handle Class 8 trucks in winter conditions. Estimated cost without any
practical alternate: \$2.4B (current) + \$0.9B (stranded freight at tail events where
I-40 is also affected by weather) = \$3.3B/year.

\textbf{With a hypothetical I-70W alternate} (western Colorado mountain corridor extended
to Sacramento, representing an upgraded Trans-Sierra route): the detour from Donner to
Sacramento via a parallel Sierra crossing would be approximately 90 miles (compared to
550 miles via I-40). B1 penalty multiplier drops from 6.50 to $1 + 90/100 = 1.90$.
Rerouting cost drops from \$2,476/truck to \$450/truck. Recalculating annual cost at
the reduced rerouting penalty: $50 \times 8,000 \times [0.38 \times \overline{C}_\text{wait}
+ 0.62 \times \$450] = 50 \times 8,000 \times (466.6 + 279) = 50 \times 8,000 \times
\$745.6 = \$298.2\text{M} \approx \$0.5\text{B}$ (including tail correction).

\textbf{Redundancy value} = current cost (\$2.4B) minus cost with I-70W alternate
(\$0.5B) = \textbf{\$1.9 billion per year}. This is the upper-bound economic benefit
of constructing an interstate-quality alternate Sierra Nevada crossing.

\subsection{Dallas Interchange: Redundancy Value \$0.4 Billion per Year}

Current: I-35E/I-35W/I-45 convergence with I-20 as the primary alternate (200-mile
detour south through Waco). Proposed alternate designation: upgrade US-287 (Dallas to
Fort Worth to Wichita Falls) to interstate standard, reducing the effective detour
to 80 miles.

Rerouting cost at 200-mile detour (current): $\$225 \times (200/60) = \$750$/truck.
Rerouting cost at 80-mile detour (US-287 upgraded): $\$225 \times (80/60) = \$300$/truck.
Closure cost reduction: $6 \times 18,000 \times 0.62 \times (\$750 - \$300) = 6 \times
18,000 \times 0.62 \times \$450 = \$301\text{M}$, rounded to \textbf{\$0.4B/yr}
accounting for partial tail-event benefits.

\subsection{General Principles of Redundancy Value}

The redundancy value analysis across all 15 corridors reveals two principles:

\textbf{Principle 1}: Redundancy value is highest when both B1 score (isolation) and
closure frequency are high. Donner (B1 score 9.2, 50 closures/yr) has the highest
redundancy value. I-95 Baltimore (B1 score 3.1, 20 closures/yr) has low redundancy
value because adequate alternates already exist (I-695 bypass): redundancy value
approximately \$0.15B/yr.

\textbf{Principle 2}: Redundancy investment and hardening investment are additive, not
substitutable, for compound exposure corridors. For Donner Pass:
\begin{itemize}
\item Hardening alone (snowsheds over the 14-mile Donner Summit section): reduces closure
frequency from 50/yr to 30/yr; reduces annual cost from \$2.4B to \$1.4B; saves \$1.0B/yr.
\item Redundancy alone (I-70W alternate): reduces per-closure rerouting cost; saves
\$1.9B/yr.
\item Combined (hardening + redundancy): closure frequency 30/yr, rerouting cost at I-70W
levels; annual cost drops to \$0.3B/yr; combined savings = \$2.1B/yr.
\end{itemize}
The combined savings (\$2.1B) exceed either investment alone because they operate on
different parameters of the cost function: hardening reduces $\lambda$ (frequency),
redundancy reduces $C_\text{reroute}$ (rerouting cost). They are complementary investments.

\begin{table}[ht]
\centering
\caption{Redundancy value for top-10 closure-cost corridors. Current cost (from
Table~\ref{tab:closure-costs}), cost with best practical alternate, and redundancy
value (\$B/yr). Best practical alternate defined as the alternate that can be built
or upgraded to interstate quality within ROUTE v1.2 scope.}
\label{tab:redundancy-value}
\begin{tabular}{@{}llcccc@{}}
\toprule
Rank & Corridor & Current & Cost w/ & Redundancy & Alternate \\
     &          & Cost    & Alternate & Value     & Description \\
     &          & (\$B/yr) & (\$B/yr) & (\$B/yr) & \\
\midrule
1 & I-80 Donner          & 2.40 & 0.50 & 1.90 & Hypothetical I-70W \\
2 & I-10 Gulf Coast LA   & 1.20 & 0.55 & 0.65 & I-20 via Shreveport \\
3 & Dallas Interchange   & 0.80 & 0.40 & 0.40 & US-287 to interstate \\
4 & I-35 Oklahoma        & 0.40 & 0.25 & 0.15 & I-44/US-77 upgrade \\
5 & I-90 Snoqualmie      & 0.31 & 0.08 & 0.23 & US-2 Stevens Pass \\
6 & I-5 Siskiyou         & 0.24 & 0.10 & 0.14 & US-97 to interstate \\
7 & I-70 Vail/Eisenhower & 0.20 & 0.09 & 0.11 & US-40/SH-9 upgrade \\
8 & I-90 Homestake MT    & 0.16 & 0.10 & 0.06 & US-12 upgrade \\
9 & I-95 Baltimore       & 0.60 & 0.45 & 0.15 & I-695 (exists) \\
10 & I-10 Gulf Coast TX  & 0.27 & 0.12 & 0.15 & I-20 via Midland \\
\bottomrule
\end{tabular}
\end{table}

\subsection{Implication for Investment Strategy}

Corridors with high redundancy value justify \textit{both} alternate route investment
\textit{and} hardening investment simultaneously, because the two address different
failure modes and their benefits are additive. The compound exposure investment logic
\citep{ROUTE_B3} follows directly: a corridor that scores high on D1 (climate exposure)
and high on B1 (isolation) is a compound risk, and investment should target both
dimensions together.

Donner Pass is the clearest case. A snowshed program that reduces closure frequency
(addressing D1) and a parallel-corridor designation or construction (addressing B1)
together produce \$2.1B/yr in benefits. Either alone captures less than half that total.
The correct investment frame for Donner is not ``hardening vs.\ redundancy'' but
``hardening and redundancy, sequenced to maximize early benefit.''
